\section{Method Objective and Design Principles}

The methodology targets sparse-demonstration humanoid object relocation, where valid solutions are multi-modal and long-horizon execution requires both local realism and task consistency. The design follows three principles:
\begin{enumerate}
    \item \textbf{Invariance by representation:} encode trajectories in robot-centric coordinates to reduce sensitivity to global pose changes.
    \item \textbf{Distributional trajectory modeling:} learn a conditional trajectory distribution instead of a single deterministic predictor.
    \item \textbf{Composable generation:} synthesize new trajectories by stitching short windows and composing part-wise motion patterns learned across different demonstrations.
\end{enumerate}

\section{Conditional Diffusion for Windowed and Long-Horizon Planning}

Let $x_0\in\mathbb{R}^{T\times D}$ be a clean trajectory window and $c=(h_t,g)$ denote history and task conditions. The model learns
\[
p_{\theta}(x_0\mid c)
\]
via a discrete diffusion process with forward corruption
\[
q(x_k\mid x_{k-1})=\mathcal{N}(\sqrt{1-\beta_k}x_{k-1},\beta_k I), \quad k=1,\dots,K,
\]
and closed form
\[
x_k=\sqrt{\bar\alpha_k}x_0+\sqrt{1-\bar\alpha_k}\,\epsilon,\qquad \epsilon\sim\mathcal{N}(0,I).
\]

The denoiser is trained with $v$-prediction:
\[
\mathcal{L}_{\mathrm{diff}}=\mathbb{E}_{x_0,\epsilon,k}\left[w(k)\,\|v-v_{\theta}(x_k,k,c)\|_2^2\right],
\qquad
v=\sqrt{\bar\alpha_k}\,\epsilon-\sqrt{1-\bar\alpha_k}\,x_0.
\]

To produce long trajectories, inference rolls out autoregressively over windows:
\[
\tau^{(m)}=\mathcal{G}_{\theta}(h^{(m)},g),
\qquad
h^{(m+1)}=\operatorname{Tail}(\tau^{(m)}),
\quad m=1,\dots,M.
\]
Hence long-horizon motion is generated by composing local segments learned from many demonstrations, including part-wise combinations (e.g., lower-body locomotion with upper-body manipulation phases) that were not seen as one complete sequence in the dataset.

\section{Trajectory Representation and Conditioning}

\subsection{Egocentric feature representation}

Each frame is represented by semantically grouped kinematic channels (reference $D=51$):
\[
[\Delta xy,\Delta\psi,\Delta xy_{obj},z_{obj},q_{joints},z_{body},r^{6D}_{body},p^{rel}_{obj},r^{6D,rel}_{obj}].
\]
This representation keeps locomotion, articulation, and object interaction in one sequence while reducing sparsity through egocentric alignment.

Formally, if $p_t^{w}$ is a world-frame planar position and $(p_0^{w},\psi_0)$ is the anchor pose, the local-frame coordinate is
\[
p_t^{loc}=R_z(-\psi_0)\,(p_t^{w}-p_0^{w}).
\]
At reconstruction time, the inverse transform is
\[
p_t^{w}=R_z(\psi_0)\,p_t^{loc}+p_0^{w}.
\]
This forward/inverse mapping is central to both training-data construction and world-frame rollout reconstruction.

\subsection{Dual conditioning, state in-painting, and indicator awareness}

State history and goal conditions are embedded separately and fused in the denoiser. During training, classifier-free condition dropout supports robust conditional generation and inference-time guidance:
\[
\hat y = y_{\mathrm{uncond}} + w_{\mathrm{cfg}}\big(y_{\mathrm{cond}}-y_{\mathrm{uncond}}\big).
\]

For masked in-painting, constrained channels are modeled as masked completion. Let $M\in\{0,1\}^{T\times D}$ be a feature-time indicator mask and $W_k$ known values at denoising step $k$:
\[
\tilde x_k=M\odot W_k + (1-M)\odot x_k.
\]
The indicator mask is embedded as an explicit conditioning signal so the model distinguishes known and unknown channels, including partial per-feature constraints.

To reduce autoregressive drift, the generated future window includes an overlap with known history frames, and those overlap states are in-painted as hard observations. In effect, the model denoises around fixed recent states instead of freely re-synthesizing them. This improves boundary continuity across stitched windows and reduces accumulation of joint-limit and posture artifacts over long horizons.

RePaint-style resampling is used to improve boundary consistency under hard constraints. For resampling iteration $r$:
\[
x_k^{(r+1)}
=
M\odot W_k
+(1-M)\odot\left(\sqrt{\bar\alpha_k}\,\hat x_0^{(r)}+\sqrt{1-\bar\alpha_k}\,\xi^{(r)}\right),
\qquad
\xi^{(r)}\sim\mathcal{N}(0,I),
\]
followed by denoising back toward lower noise levels while reapplying $M$. This iterative re-noise/re-denoise cycle reduces seam artifacts between constrained and unconstrained channels.

\subsection{Data augmentation and preprocessing operations}

Data operations are treated as part of the method, not only as implementation details, as it helps increase the coverage of the data over the feature space (or reduces sparsity of data).

\paragraph{Mirror symmetry augmentation.}
The dataset applies left-right sagittal reflection with probability $p_{\mathrm{mirror}}$, including feature permutation and sign transforms that map each left-body channel to its right-body counterpart (and vice versa), while preserving physically meaningful coordinate conventions. Base-orientation, lateral displacement, and object-relative channels are transformed consistently under the same reflection map. This effectively doubles directional skill coverage from sparse demonstrations, reduces lateral bias, and improves symmetry-consistent generalization under mirrored goals.

\paragraph{Goal-vector normalization and magnitude clipping.}
Given local goal displacement $\Delta g\in\mathbb{R}^{3}$, the task vector is encoded as direction plus clipped magnitude:
\[
d=\frac{\Delta g}{\|\Delta g\|_2+\varepsilon},
\qquad
m=\min\big(\|\Delta g\|_2,\,m_{\max}\big),
\qquad
\tilde g=[d,m].
\]
The separation of direction and magnitude stabilizes conditioning across large displacement ranges and allows learning motion priors for a general direction from a lot of motions. The magnitude clipping allows the model to view out-of-distribution, larger distances as within distribution. This works particularly when the clipping distance is picked such that the robot would have to turn or walk in order to place the object at the desired distance, as it would have to do the same action for a longer goal in the same direction.

\paragraph{Stochastic conditioning perturbations.}
Optional observation and goal noise are injected during training to reduce overfitting to exact conditioning vectors and improve robustness under rollout mismatch. The noise levels are tuned per feature group based on how much error (from sensors in hardware, or motion generation in simulation) is to be considered for the model predictions. 

\subsection{Ablation-informed design choices}

The final method configuration is guided by ablations under sparse-data training (four reference motions, densified through trajectory optimization and augmentation). The following practical choices were retained:
\begin{itemize}
    \item \textbf{Backbone selection:} the transformer denoiser is preferred over a U-Net alternative for better modeling of higher-frequency temporal variations in joint-object interaction windows.
    \item \textbf{Optimization scale:} a large batch regime (reference value: 1024) gives stable gradient estimates while maintaining practical update frequency per epoch.
    \item \textbf{Classifier-free conditioning:} condition dropout improves controllability through guidance, with the known trade-off that overly strong guidance on highly out-of-distribution tasks can reduce physical plausibility.
    \item \textbf{Goal encoding strategy:} direction-normalized and magnitude-clipped goal encoding improves extrapolation to longer relocation distances by reusing turning/walking primitives instead of overfitting to narrow training distances.
    \item \textbf{Rotation parameterization:} a 6D orientation representation is used to avoid discontinuities common in Euler-angle or quaternion parameterizations when learned directly in sequence models.
    \item \textbf{State-conditioning masking:} random masking of conditioning feature groups reduces spurious history-future coupling and encourages reliance on semantically relevant cues rather than shortcut correlations.
    \item \textbf{Redundant object channels for control:} explicit object displacement channels ($\Delta xy_{obj}, z_{obj}$) are retained to support constrained waypoint in-painting, improving target convergence while accepting a potential trade-off in strict physical consistency.
\end{itemize}
These observations motivate the masking and constrained-completion formulation in the next section.

\section{Masking, Motion In-Betweening, and Motion In-Painting}

This work treats trajectory generation as a constrained completion problem in addition to standard denoising. Instead of always predicting all channels at all timesteps, the model is trained with structured masks that define known versus unknown states, and learns to complete missing motion consistent with both dynamics priors and task conditions.

\subsection{Feature-time masking formulation}

Let $M\in\{0,1\}^{T\times D}$ be a binary mask over time and feature channels, where $M_{t,d}=1$ marks known values. At diffusion step $k$, the model input is constrained as
\[
\tilde x_k = M\odot W_k + (1-M)\odot x_k,
\]
where $W_k$ denotes known waypoint values (clean or noise-matched depending on regime). A separate mask-indicator embedding is injected into the denoiser so that the network is explicitly aware of which channels are constrained.

In addition to waypoint masks over trajectory tokens, state-conditioning masks are sampled on the history input. Let $c_h$ be the history-condition vector and $m_h\in\{0,1\}^{d_h}$ a random feature-group mask. The effective condition is
\[
\tilde c_h = m_h\odot c_h.
\]
This prevents the denoiser from over-relying on fixed history-future correlations and encourages robust feature-wise credit assignment (e.g., locomotion cues versus manipulation cues).

To support constraint-driven generation, redundant object-motion channels (notably $\Delta xy_{obj}$ and $z_{obj}$) are retained in the trajectory representation. These channels enable direct waypoint specification for object motion while the model infers remaining degrees of freedom (base pose, joints, orientation) to remain consistent with the supplied object path.

\subsection{Motion in-betweening during training}

In-betweening is implemented by sampling sparse keyframes and requiring the network to reconstruct intermediate motion. The masking policy includes:
\begin{itemize}
    \item \textbf{Full keyframes:} all channels at selected timesteps are known.
    \item \textbf{Partial keyframes:} only selected semantic groups (e.g., locomotion, upper/lower body, object pose) are known.
    \item \textbf{Extreme conditioning cases:} occasional all-known windows and near-empty windows (single group at one frame) to improve robustness.
\end{itemize}

For fully known keyframes, the corresponding diffusion noise level is set to $k=0$ so these anchors are treated as clean constraints. For partially known keyframes, known channels are injected either as clean values or as noise-matched values at the current $k$, enforcing consistency between constrained and unconstrained channels across denoising steps.

This transforms training from pure next-trajectory synthesis into \emph{constrained trajectory completion}, improving interpolation quality and temporal coherence between sparse waypoints.

The key design is train-inference symmetry: during training, random semantic subsets are revealed and the model learns to complete the complement; during inference, planned object waypoints (from heuristics or a separate predictor) can be revealed in the same format. The denoiser then performs conditional completion under hard object-trajectory constraints. Empirically, this improves end-position convergence of the manipulated object, with a known trade-off that strict waypoint forcing can reduce global physical naturalness if waypoints are poorly chosen.

\subsection{Motion in-painting for autoregressive rollout}

During long-horizon rollout, each generated window overlaps with recent history. In-painting pins overlap states to known observations, so the model denoises around fixed boundary conditions instead of re-generating the boundary freely. Conceptually, for overlap index set $\Omega$,
\[
\hat x_{t,d} \leftarrow x^{\mathrm{obs}}_{t,d}, \qquad (t,d)\in\Omega.
\]

This reduces seam artifacts between consecutive windows, limits drift accumulation, and improves state continuity (especially joint pose, base orientation, and object-relative channels).

To further stabilize hard constraints, a RePaint-style resampling loop is used: re-noise to a higher diffusion level, re-impose constraints, and denoise again. Repeating this cycle harmonizes constrained and unconstrained dimensions near boundaries and improves stitched long-horizon consistency.

\section{Network Architecture}

The trajectory denoiser is implemented as a conditional \textbf{DiT-1D} backbone operating on a fixed horizon of motion tokens. For each sample, the network receives a noisy trajectory tensor
\(x_k \in \mathbb{R}^{T\times C}\), where in our setup \(T=20\) and \(C=51\), and predicts the reverse-diffusion target (velocity parameterization).

\paragraph{1) Structured token embedding}
Instead of a single flat projection for all channels, the model uses \emph{group-wise semantic projections} over feature groups (root motion, object motion, joints, body orientation, etc.). For timestep \(t\):
\[
h_t=\sum_{m\in\mathcal{M}} W_m x_t^{(m)} + b_m,
\]
where \(x_t^{(m)}\) is the sub-vector of group \(m\), and each group has its own learned linear map to the common hidden dimension \(d\) (here \(d=256\)). This preserves modality-specific inductive bias while still producing one fused token per frame for global self-attention.

\paragraph{2) Positional and diffusion-step conditioning}
The embedded sequence is augmented with learnable 1D temporal positional embeddings:
\[
z_t^{(0)} = h_t + p_t.
\]
Diffusion noise levels are embedded per token and injected as conditioning. The implementation supports token-wise noise levels (including feature-wise schedules reduced to token level for conditioning), which is useful for masked/in-betweening supervision.

\paragraph{3) External condition encoder (state history + task goal)}
Conditioning is encoded in two branches and fused in latent space:
\begin{itemize}
    \item \textbf{State-history branch:} each observed state frame is mapped by an MLP to \(d\)-dimensional embeddings, then aggregated into a single context token using learned attention pooling (or mean pooling in ablations).
    \item \textbf{Task branch:} goal/task parameters are mapped by a second MLP to a \(d\)-dimensional token.
\end{itemize}
Both branches use classifier-free style dropout during training (state/task dropped stochastically), and their embeddings are concatenated then projected into the backbone conditioning space. This yields robust conditional and unconditional behavior under one shared model.

\paragraph{4) DiT block design (AdaLN-Zero modulation)}
The backbone uses \(L=8\) transformer blocks with \(4\) attention heads and MLP expansion ratio \(4\). Conditioning is applied with \textbf{AdaLN-Zero} rather than cross-attention. This is the default backbone configuration used in the thesis: it was selected as a compact and stable baseline rather than through a dedicated depth/head-count sweep. For each block:
\[
\tilde{x}=\mathrm{LN}(x)\odot(1+\gamma(c))+\beta(c),
\]
\[
x \leftarrow x + g_{\text{attn}}(c)\odot \mathrm{MSA}(\tilde{x}),
\]
\[
x \leftarrow x + g_{\text{mlp}}(c)\odot \mathrm{MLP}(\tilde{x}),
\]
where \(\beta,\gamma,g\) are predicted from condition \(c\). Zero-initialized modulation/gates make blocks near-identity at initialization, improving optimization stability for diffusion training.

\paragraph{5) Output head and prediction target}
After the transformer stack, a final adaptive normalization layer and linear projection map tokens back to \(C=51\) channels per timestep. The diffusion objective uses \(v\)-prediction with a discrete noise schedule (train-time 200 steps, accelerated inference with 10 steps), giving stable denoising and efficient rollout.

Overall, this architecture combines: (i) semantically structured input projections, (ii) strong temporal self-attention, and (iii) robust conditional modulation through AdaLN-Zero and condition dropout, which together improve trajectory realism, controllability, and generalization.

\begin{figure}[!h]
    \centering
    \includegraphics[width=0.88\linewidth]{diffusion-architecture.png}
    \caption{Network architecture used for conditional trajectory diffusion in this thesis.}
    \label{fig:method_diffusion_architecture}
\end{figure}

\section{Physics-Aware Auxiliary Losses}

Pure diffusion reconstruction can produce locally noisy or weakly consistent kinematic plans under sparse data. We therefore use auxiliary terms as soft physics-aware priors (without claiming full dynamic feasibility). These losses are applied on predicted clean trajectories and only over valid frame pairs, so padding or invalid tokens do not bias optimization.

Let $m_t\in\{0,1\}$ be a frame-validity indicator and
\[
\omega_t = m_t m_{t+1}, \qquad t=1,\dots,T-1
\]
be the pair-validity mask used for consecutive-frame losses.

\paragraph{Temporal smoothness.}
Rather than only shrinking motion magnitude, smoothness is imposed by matching predicted and reference first-order temporal differences:
\[
\mathcal{L}_{\mathrm{smooth}}=
\frac{1}{\sum_t \omega_t + \varepsilon}
\sum_{t=1}^{T-1}\omega_t\,
\left\| (\hat x_{t+1}-\hat x_t) - (x^{gt}_{t+1}-x^{gt}_t) \right\|_2^2.
\]
In practice, this term can be restricted to selected feature groups (e.g., joints and body orientation) to avoid over-regularizing channels that should remain agile.

\paragraph{Grasp-consistency.}
Let $p_t^{rel}$ denote object position in the body frame. During grasp/carry phases, this relative displacement should vary smoothly. We use a distance-based soft gate
\[
\gamma_t = \sigma\!\big(\kappa(\tau_g-\|p_t^{rel,gt}\|_2)\big),
\]
with threshold $\tau_g$ and sharpness $\kappa$, so supervision is strong when the object is close and weak otherwise. The loss is
\[
\mathcal{L}_{\mathrm{grasp}}=
\frac{1}{\sum_t\omega_t\,\bar\gamma_t+\varepsilon}
\sum_{t=1}^{T-1}
\omega_t\,\bar\gamma_t\,
\left\| (p_{t+1}^{rel}-p_t^{rel})-(p_{t+1}^{rel,gt}-p_t^{rel,gt}) \right\|_2^2,
\quad
\bar\gamma_t = \min(\gamma_t,\gamma_{t+1}).
\]

\paragraph{Optional feet-sliding proxy.}
An additional proxy can penalize ankle-joint velocity during stance-like intervals. With predicted ankle velocities $v_t^{ankle}$ and a stance gate $\eta_t$ derived from low reference ankle speed,
\[
\mathcal{L}_{\mathrm{feet}}=
\frac{1}{\sum_t\omega_t\,\eta_t+\varepsilon}
\sum_{t=1}^{T-1}\omega_t\,\eta_t\,\|v_t^{ankle}\|_2^2.
\]
This term is optional and mainly used as a regularizer against visible foot jitter in contact phases.

The total objective is
\[
\mathcal{L}=\mathcal{L}_{\mathrm{diff}}
+\lambda_s\mathcal{L}_{\mathrm{smooth}}
+\lambda_g\mathcal{L}_{\mathrm{grasp}}
+\lambda_f\mathcal{L}_{\mathrm{feet}}.
\]
In the default setting, smoothness and grasp-consistency are active, while the feet term is enabled only when needed for additional contact-phase stabilization.

